\section{验证架构、测试目录与外部对照边界}
\label{sec:scenario-validation-architecture}

\subsection{验证金字塔}
\begin{center}\small
\begin{tabular}{@{}L{3.2cm}L{5.3cm}L{5.8cm}@{}}
\toprule
层级 & 证据 & 能证明的对象\\
\midrule
L1 解析/不变量 & 独立公式、概率和、有限数 & 局部方程与守恒\\
L2 单元回归 & Catch2 固定 seed/边界案例 & 分支、fallback、字段消费\\
L3 交叉实现 & C++ driver + 独立 Python & 序列化结果可由公开方程复算\\
L4 接口/模式 & JSON v3、workbook、HTTP & 数据契约和生产路由\\
L5 下游任务 & PF/OPF/reliability/resilience 回测 & 代表场景对具体决策任务的充分性\\
\bottomrule
\end{tabular}
\end{center}
当前模块具备 L1--L4 的专项证据；L5 需要按下游模块另行设计，不能由场景生成测试代替。

\subsection{test\_scenario\_generation 测试目录}
\begin{longtable}{@{}L{8.0cm}L{6.0cm}@{}}
\toprule
Catch2 case & 验证对象\\
\midrule\endfirsthead
\toprule Catch2 case & 验证对象\\\midrule\endhead
Regular scenario generation applies SSP/year climate morphing to load and PV & 已知 SSP/year 气候形变\\
Scenario generation prepares standard multiplier time series internally & multiplier/base/capacity 契约\\
Regular climate perturbation creates reproducible candidate variation & 固定 seed 与气候随机性\\
Scenario perturbations realize declared temporal and cross correlations & AR(1) 与跨变量统计目标\\
Hybrid clustering consumes the declared quantile boundary fraction & boundary fraction 字段消费\\
Typhoon minimum fault probability gates sampling without hiding risk & sampling gate 与风险保留\\
Regular climate morphing falls back safely for unknown SSP/year & 零增量 fallback/audit\\
Large annual scenario generation uses bounded aggregate load profiles & 年度内存预算与 aggregate 降级\\
Large reliability catalog bounds audit samples without dropping contingencies & audit 预算不删目录\\
Very large reliability catalogs preserve N-1 coverage within work budgets & 大目录 N--1 覆盖\\
Typhoon fault generation reuses precomputed topology segments & 空间预计算复用\\
Large typhoon fault sets use bounded upstream-first repair ordering & 近似修复排序与标志\\
\bottomrule
\end{longtable}

\subsection{test\_typhoon\_traffic\_impact 测试目录}
\begin{longtable}{@{}L{8.0cm}L{6.0cm}@{}}
\toprule
Catch2 case & 验证对象\\
\midrule\endfirsthead
Typhoon road impact creates common speed, capacity, and closure profiles & 逐 link profile 与闭锁\\
Typhoon road impact marks affine georeferencing of benchmark topology & Auto/affine fallback 标志\\
CTM reads road profiles by parent simulation step & 下游时步映射\\
Typhoon catalog snapshots survive concurrent mixed-option refreshes & 缓存线程安全快照\\
Typhoon road impact rejects invalid clamp and hydrology bounds & 选项失败语义\\
\bottomrule
\end{longtable}

\subsection{test\_scenario\_bundle\_schema 测试目录}
\begin{longtable}{@{}L{8.0cm}L{6.0cm}@{}}
\toprule
Catch2 case & 验证对象\\
\midrule\endfirsthead
Scenario bundle validation accepts v3 dimensionless multiplier schema & v3 正常路径\\
Scenario bundle validation rejects invalid multiplier profiles & 负值/SOC/长度等错误\\
Scenario bundle normalization upgrades v2 case shape & v2→v3 迁移与 warning\\
\bottomrule
\end{longtable}

\subsection{独立 oracle 设计}
\srcpath{tests/scenario\_generation\_cross\_validation.py} 不链接也不导入 hacdcpf。它从 C++ JSON
读取公开输入，独立实现 Haversine、Holland、降雨、等效风、架空段失效、交通积水/容量、
Pearson 统计、median/IQR 缩放和簇距离复算。验证项包括：
\begin{itemize}
  \item 候选和代表概率误差 $<10^{-12}$；
  \item 4096 步 load lag-1 与目标差 $<0.05$；全天 load-wind 同期相关与目标差 $<0.08$；
  \item 风、雨、易损概率和交通递推误差 $<10^{-10}$；
  \item 128 个成员无遗漏、代表属于成员、平均/最大距离和 transport 指标误差 $<10^{-10}$；
  \item sampling gate 下 fault count 为 0 且风险概率仍保留。
\end{itemize}

\subsection{为什么 cross-correlation 使用风电}
PV 基线夜间为零，只能在白天恢复 multiplier。第一次预声明容差为 0.08 时，PV 子样本实测
-0.341279，相对目标 -0.25 偏差 0.091279，验证失败。没有放宽阈值，而是重新推导实验：使用
全天非零 wind profile 恢复完整 renewable multiplier，实测 -0.290252，偏差 0.040252，通过
原阈值；PV 值继续报告为周期性子样本诊断。这是测试设计修复，不是结果筛选。

\subsection{OpenDSS 与 GridLAB-D 的正确角色}
OpenDSS 和 GridLAB-D 能消费确定性负荷/发电 profile 并运行时序潮流，但它们没有与本模块
等价的 SSP 气候形变、台风目录、易损故障抽样和尾部锚定 k-medoids API。因此：
\begin{itemize}
  \item 可以把代表 profile 导出后比较下游电压/潮流；
  \item 不能把下游潮流一致写成“场景生成算法认证”；
  \item 本手册没有伪造两软件的聚类或灾害链数值；
  \item 若未来增加下游 profile replay，应单独记录软件版本、网络转换、控制模式和误差。
\end{itemize}

\subsection{统计准入与防止移动门槛}
容差必须在运行前声明。若实测偏差超限，按以下顺序处理：
\begin{enumerate}
  \item 检查 oracle 是否独立、输入是否相同、单位/索引是否一致；
  \item 检查有限样本、AR 有效样本量、clip 和周期性选择；
  \item 检查 C++ 是否忠实实现理论；
  \item 只有理论重新推导支持时才修改实验或门槛，并保留失败原因。
\end{enumerate}

\subsection{未覆盖测试}
当前尚无：历史台风事件的空间观测评分、校准易损曲线的 out-of-sample 检验、场景数序贯收敛、
下游 OPF 决策 regret、跨平台 bitwise 重现、v3 工作簿大规模 round-trip 性能基线。这些必须
列为未认证边界，不能写入“已全面验证”。
